\documentclass[main.tex]{subfiles}
\begin{document}
\ifSubfilesClassLoaded{\begin{center}\small\itshape Part of the GnssProcessor Mathematical Reference \textbullet\ \today\end{center}}{}
\section{Site Displacement and Geophysical Corrections}\label{sec:sitedisp}
% ============================================================
%  Geophysical station-position displacements applied to the
%  receiver coordinate BEFORE the geometric range is formed.
% ============================================================

% --- Doc-local symbols ---------------------------------------------------
\newcommand{\disp}{\bm{d}}                       % 3D ECEF site displacement (m)
\newcommand{\dispSET}{\disp_{\mathrm{SET}}}      % solid Earth tide displacement
\newcommand{\dispOTL}{\disp_{\mathrm{OTL}}}      % ocean tidal loading displacement
\newcommand{\dispPT}{\disp_{\mathrm{PT}}}        % (solid) pole tide displacement
\newcommand{\dispOPT}{\disp_{\mathrm{OPT}}}      % ocean pole tide displacement
\newcommand{\dispAPL}{\disp_{\mathrm{APL}}}      % atmospheric pressure loading
% \leftmoon is not provided by the math packages loaded in gnssreport.sty
% (amsmath/amssymb/amsthm/mathtools/bm); define a doc-local Moon glyph so the
% Sun/Moon body-tide notation is self-contained without wasysym/marvosym.
\newcommand{\leftmoon}{\mathrm{m}}               % Moon symbol (paired with \odot)
\newcommand{\rxsun}{\bm{r}^{\odot}}              % geocentric Sun position
\newcommand{\rxmoon}{\bm{r}^{\leftmoon}}         % geocentric Moon position
\newcommand{\loveh}{h}                              % degree-2 Love number
\newcommand{\shidal}{l}                              % degree-2 Shida number
\newcommand{\geolat}{\varphi}                    % (geocentric or geodetic) latitude
\newcommand{\geolon}{\lambda}                    % longitude
\newcommand{\unitup}{\hat{\bm{u}}}               % local radial / up unit vector
\newcommand{\polem}[1]{m_{#1}}                   % polar-motion wobble m1,m2

The receiver coordinate that enters \dcref{eq:los}{the line-of-sight
construction} and \dcref{eq:geom-linear}{the linearised geometric range}
of the Foundations document is not the instantaneous position of the
crust.  It is the conventional tide-free station position, to which the
IGS/ITRF products are referenced, plus a sum of \emph{a~priori}
geophysical corrections.  Each correction is a modelled 3D displacement
$\disp(t)$ of the antenna reference point.  The ARP is the position
state, which the engine obtains from the RINEX monument and the header
eccentricity (\dcref{sec:frame-arp}{the marker/ARP section of the
Coordinate Frames document}).  None of the corrections is estimated.
All are computed from the a~priori position (good to centimetres), the
Sun and Moon ephemerides, and the Earth-orientation parameters.  They
are then projected onto the line of sight and added to the modelled
range.  This document derives the models as the IERS~2010 conventions
\cite{iers2010} intend them, so it can be used to settle questions of
sign and phase in the implementation.

Write the corrected receiver position as
\begin{equation}
\rxpos = \rxpos_0 + \dispSET + \dispOTL + \dispPT + \dispOPT + \dispAPL ,
\label{eq:sitedisp-sum}
\end{equation}
with $\rxpos_0$ the conventional a~priori position and each $\disp$ an
ECEF vector of millimetres to decimetres.  The five terms are:
\begin{itemize}
\item the solid Earth tide $\dispSET$, the largest, up to $\sim\!30$~cm
      radially;
\item ocean tidal loading $\dispOTL$, centimetres near coasts;
\item the solid (rotational-deformation) pole tide $\dispPT$, below a
      centimetre to a centimetre, secular plus Chandler;
\item the ocean pole tide $\dispOPT$, below a millimetre to
      $\sim\!2$~mm;
\item atmospheric pressure loading $\dispAPL$, $2$--$20$~mm.
\end{itemize}
Only the sun-synchronous $S_1/S_2$ thermal-tide part of the atmospheric
loading is modelled.  The non-tidal (synoptic) part is a known gap
(\Cref{sec:sitedisp-apl}).  All five are computed in
\texttt{src/core/corrections/GnssCorrections.cpp} and
\texttt{src/core/corrections/OceanTideModel.cpp}.  The LOS projection
and the summation into the modelled range are in
\texttt{src/core/model/ModelGeometry.cpp} and
\texttt{src/core/model/PppModel.cpp}.  Each term has a config flag
(\texttt{enable\_solid\_earth\_tide}, \texttt{enable\_ocean\_tide\_loading},
\texttt{enable\_pole\_tide}, \texttt{enable\_ocean\_pole\_tide},
\texttt{enable\_atmospheric\_loading}).  With a flag off, that term is
a zero vector.

\subsection{Sign convention and projection into the range}
\label{sec:sitedisp-projection}

Every displacement is projected onto the unit line of sight before it
enters the modelled range.  Let
$\los = (\satpos - \rxpos)/\lVert\satpos-\rxpos\rVert$ be the
receiver-to-satellite unit vector (the convention of
\dcref{eq:los}{the Foundations LOS equation}; see
\texttt{ModelGeometry.cpp} where \texttt{los = (sat\_pos - rx\_pos)/range}).
A site displacement $\disp$ changes the geometric range by its component
along the line of sight.  The range shortens when the antenna moves
\emph{toward} the satellite.  The scalar correction added to the
modelled range is
\begin{empheq}[box=\fbox]{equation}
\delta\rho \;=\; \disp\cdot(-\los) \;=\; \disp\cdot\frac{\rxpos-\satpos}{\lVert\rxpos-\satpos\rVert},
\label{eq:sitedisp-projection}
\end{empheq}
so a displacement toward the satellite (a positive component along
$\los$) gives $\delta\rho<0$.  The ocean-loading, atmospheric-loading,
and pole-tide projections use this form
(\texttt{Geometry::computeSatelliteGeometry}: \texttt{otl\_corr = otl\_ecef.dot(-los)},
\texttt{apl\_corr = apl\_ecef.dot(-los)};
\texttt{PppModel::computePoleTideProjection}: \texttt{pole\_tide\_ecef.dot(-los)}), and the
modelled range adds them in (\texttt{PppModel::computeBaseModeledValue}
return sums \texttt{+ otl\_corr + apl\_corr + solid\_tide} and the pair
processors add \texttt{+ pole\_tide\_corr}).
The solid-Earth-tide helper does its own projection with the same vector
$(\rxpos-\satpos)/\lVert\rxpos-\satpos\rVert$
(\texttt{GnssCorrections.cpp}, \texttt{computeSolidEarthTide}:
\texttt{los = (receiver\_pos - sat\_pos)/rho\_norm}, then
\texttt{disp.dot(los)}), so all the tide terms use one sign convention.
The rest of this document gives each model as a 3D ECEF (or local)
displacement $\disp$ and leaves \Cref{eq:sitedisp-projection} implicit.

The Sun and Moon positions $\rxsun,\rxmoon$ on which the body tides
depend, and the ECI$\leftrightarrow$ECEF rotation and Earth-orientation
parameters they require, come from
\texttt{src/core/model/Planetary.cpp} and are documented in
\dcref{sec:frames}{the Coordinate Frames and Time Systems document}.
Here $\rxsun,\rxmoon$ are given geocentric ECEF vectors.

% ============================================================
\subsection{Solid Earth tide (IERS 2010, Steps 1 + 2)}
\label{sec:sitedisp-set}

The solid Earth tide is the elastic deformation of the whole Earth under
the lunisolar tide-generating potential.  The implementation is a port
of the IERS Conventions \texttt{dehanttideinel} routine \cite{iers2010}
and has two steps.  Step~1 is evaluated in the time domain from the
instantaneous Sun and Moon geometry.  Step~2 is a frequency-domain
correction for the few constituents whose Love and Shida numbers differ
from the nominal degree-2 values.  \textbf{The whole Step~1+2 block is
validated against reference values}: it reproduces the official IERS
\texttt{dehanttideinel} test vector and the in-tree POINT reference to
below a millimetre, cross-checked by
\texttt{scripts/solid\_tide\_reference.py}
(\texttt{GnssCorrections.cpp}, \texttt{computeSolidEarthTideDisplacement} and
the \texttt{SolidEarthTideTest} suite).  The permanent (zero-frequency)
tide is \emph{retained}, so the corrected coordinate is conventional
tide-free, as in IGS/ITRF (see the routine's header comment).

\subsubsection{Step 1: in-phase degree 2 and 3}

For each tide-raising body $b\in\{\odot,\leftmoon\}$ at geocentric
position $\bm{r}^b$ ($r_b=\lVert\bm{r}^b\rVert$) and the station at
$\rxpos_0$ ($r=\lVert\rxpos_0\rVert$, geocentric latitude $\geolat$ via
direction cosines), the degree-2 in-phase displacement is built from the
cosine of the body zenith angle
$\cos Z_b = \hat{\rxpos}_0\cdot\hat{\bm{r}}^b$ and a body scale factor
\begin{equation}
\mathrm{fac}_b \;=\; \frac{M_b}{M_\oplus}\,R_e\left(\frac{R_e}{r_b}\right)^{3},
\label{eq:sitedisp-set-fac}
\end{equation}
with $R_e = 6378136.55$~m and mass ratios $M_\odot/M_\oplus = 332945.943062$,
$M_\leftmoon/M_\oplus = 0.012300034$
(\texttt{GnssCorrections.cpp} tide-helper constants
\texttt{kTideRe}/\texttt{kTideMassRatioSun}/\texttt{kTideMassRatioMoon}).  The radial and
tangential weights are
\begin{align}
p_2^b &= 3\!\left(\tfrac{\loveh_2}{2} - \shidal_2\right)\cos^2 Z_b - \tfrac{\loveh_2}{2},
\label{eq:sitedisp-set-p2}\\
x_2^b &= 3\,\shidal_2 \cos Z_b ,
\label{eq:sitedisp-set-x2}
\end{align}
(\texttt{GnssCorrections.cpp}, \texttt{computeSolidEarthTideDisplacement}), where $\loveh_2,\shidal_2$ are the
\emph{latitude-dependent} degree-2 Love and Shida numbers of IERS
Eq.~7.2,
\begin{equation}
\loveh_2 = 0.6078 - 0.0006\,(1 - \tfrac{3}{2}\cos^2\!\geolat),\qquad
\shidal_2 = 0.0847 + 0.0002\,(1 - \tfrac{3}{2}\cos^2\!\geolat)
\label{eq:sitedisp-set-love}
\end{equation}
(\texttt{GnssCorrections.cpp}, \texttt{computeSolidEarthTideDisplacement}).  An analogous degree-3 term
uses fixed $\loveh_3=0.292,\ \shidal_3=0.015$ (\texttt{kTideH3}/\texttt{kTideL3}) and
$\mathrm{fac}_b^{(3)} = \mathrm{fac}_b\,(R_e/r_b)$
(same routine).  The Step-1
in-phase displacement sums the radial weight along the station unit
vector and the tangential weight along each body's unit vector:
\begin{empheq}[box=\fbox]{equation}
\dispSET^{(1)} = \sum_{b\in\{\odot,\leftmoon\}}\!\mathrm{fac}_b\!\left(
   x_2^b\,\hat{\bm{r}}^b + p_2^b\,\hat{\rxpos}_0\right)
   \;+\; (\text{degree-3 analogue}),
\label{eq:sitedisp-set-step1}
\end{empheq}
matching the per-component accumulation in
\texttt{computeSolidEarthTideDisplacement}.

\subsubsection{Step 1: out-of-phase anelasticity and the $L^1$ term}

Mantle anelasticity makes the diurnal and semidiurnal degree-2 Love
numbers complex; the small imaginary parts produce out-of-phase
displacements \texttt{ST1IDIU} (diurnal, $\delta\loveh=-0.0025$,
$\delta\shidal=-0.0007$) and \texttt{ST1ISEM} (semidiurnal,
$\delta\loveh=-0.0022$, $\delta\shidal=-0.0007$)
(\texttt{GnssCorrections.cpp}, helpers \texttt{st1Diurnal}, \texttt{st1SemiD}).  The latitude ($L^1$)
dependence of the Love numbers contributes a horizontal-only correction
\texttt{ST1L1} with $l^{(1)}_{\mathrm{diu}}=0.0012$,
$l^{(1)}_{\mathrm{sem}}=0.0024$
(helper \texttt{st1L1}).  These three helpers
return local radial/north/east
displacements rotated into ECEF and are added to
\Cref{eq:sitedisp-set-step1} inside \texttt{computeSolidEarthTideDisplacement}.
Their trigonometric forms follow the IERS reference and are not
reproduced here.

\subsubsection{Step 2: frequency dependence (Doodson tables)}

Step~2 corrects the constituents whose actual Love numbers differ from
the nominal degree-2 value, summing harmonic terms keyed by the Doodson
fundamental arguments.  The astronomical arguments
$\tau,s,h,p,N',p_s$ (lunar time, Moon/Sun mean longitudes, lunar
perigee, negative node, solar perigee) are polynomials in Julian
centuries $T$ from J2000, with the $15^\circ/\mathrm{hr}$ Earth-rotation
term entering $\tau$ through the fractional UT hour $H$:
\begin{align}
s   &= 218.31664563^\circ + 481267.88194\,T - 0.0014663889\,T^2 + \ldots,
\label{eq:sitedisp-doodson-s}\\
\tau &= 15^\circ\! H + 280.4606184^\circ + 36000.7700536\,T + \ldots - s,
\label{eq:sitedisp-doodson-tau}
\end{align}
with the remaining $h,p,N',p_s$ polynomials given to fourth order in
\texttt{GnssCorrections.cpp} (\texttt{computeAngles}).  For
each table line the phase is
$\theta = \tau + n_s s + n_h h + n_p p + n_{N'}N' + n_{p_s}p_s
(+\,\geolon$ for diurnal$)$ and the in-/out-of-phase radial and
tangential amplitudes (millimetres) are applied with a latitude weight.

\paragraph{Diurnal band (Table 7.3a).}  Thirty-one lines, dominated by
$K_1$ and $P_1$ (\texttt{GnssCorrections.cpp}, \texttt{st2Diurnal}).  The
displacement for one line is
\begin{equation}
\begin{aligned}
\delta R &= (A^{\mathrm{in}}_R\sin\theta + A^{\mathrm{out}}_R\cos\theta)\,
            2\sin\geolat\cos\geolat, \\
\delta N &= (A^{\mathrm{in}}_T\sin\theta + A^{\mathrm{out}}_T\cos\theta)\,
            \cos 2\geolat, \\
\delta E &= (A^{\mathrm{in}}_T\cos\theta + A^{\mathrm{out}}_T\sin\theta)\,
            \sin\geolat,
\end{aligned}
\label{eq:sitedisp-step2diu}
\end{equation}
(\texttt{st2Diurnal}), summed over the table and
rotated radial/north/east $\to$ ECEF.  The reference amplitudes (mm)
include $K_1: (12.00,\,-0.80,\,-0.67,\,-0.03)$ and
$P_1: (-1.23,\,-0.07,\,0.06,\,0.01)$ in the order
$(A^{\mathrm{in}}_R, A^{\mathrm{out}}_R, A^{\mathrm{in}}_T, A^{\mathrm{out}}_T)$
(the \texttt{st2Diurnal} table \texttt{T}).

\paragraph{Long-period band (Table 7.3b).}  Five zonal lines, with no
longitude or $\tau$ dependence and no east component
(\texttt{GnssCorrections.cpp}, \texttt{st2Lon}).  The radial weight is the
zonal Legendre form $(3\sin^2\!\geolat - 1)/2$ and the north weight is
$\sin 2\geolat$:
\begin{equation}
\delta R = (A^{\mathrm{in}}_R\cos\theta + A^{\mathrm{out}}_R\sin\theta)\,
           \frac{3\sin^2\!\geolat - 1}{2},\qquad
\delta N = (A^{\mathrm{in}}_T\cos\theta + A^{\mathrm{out}}_T\sin\theta)\,
           \sin 2\geolat,
\label{eq:sitedisp-step2lon}
\end{equation}
(\texttt{st2Lon}).  The full solid-tide
displacement is
$\dispSET = \dispSET^{(1)} + \mathrm{ST1IDIU} + \mathrm{ST1ISEM}
+ \mathrm{ST1L1} + \mathrm{STEP2DIU} + \mathrm{STEP2LON}$
(\texttt{computeSolidEarthTideDisplacement}).

% ============================================================
\subsection{Ocean tidal loading (OTL)}
\label{sec:sitedisp-otl}

Ocean tidal loading is the elastic response of the crust to the moving
mass of the ocean tides.  It reaches centimetres near coasts and is
computed as a harmonic sum over the eleven principal
constituents M2, S2, N2, K2, K1, O1, P1, Q1, Mf, Mm, Ssa
(\texttt{OceanTideModel.hpp}, \texttt{OTL\_CONSTITUENT\_NAMES}).  Each
constituent $k$ contributes, on each local component
$c\in\{\text{up},\text{west},\text{south}\}$,
\begin{empheq}[box=\fbox]{equation}
dr_c(t) \;=\; \sum_{k=1}^{11} A_{c,k}\,\cos\!\big(\chi_k(t) - \phi_{c,k}\big),
\label{eq:sitedisp-otl}
\end{empheq}
where $A_{c,k}$ is the station amplitude, $\phi_{c,k}$ the Greenwich
phase lag, and $\chi_k(t)$ the astronomical argument (the
eleven-constituent accumulation of
\texttt{OceanTideModel::computeDisplacement}).  The internal
accumulation is in the $(\text{up},\text{west},\text{south})$ frame; the
conversion to ENU is
$(E,N,U)=(-\text{west},-\text{south},\text{up})$, then ENU$\to$ECEF
(same routine).

\subsubsection{Astronomical argument and the Doodson--Warburg phase bias}
\label{sec:sitedisp-otl-dw}

The argument $\chi_k$ is assembled from Doodson multipliers
$[\,\tau,s,h,p,N',p_s\,]$ acting on the variables
$[\,\text{GMST}+\pi,\,s,\,h,\,p,\,N',\,p_s\,]$, themselves derived from
the Delaunay fundamental arguments
$s=F+\Omega,\ h=s-D,\ p=s-l,\ N'=-\Omega$
(\texttt{OceanTideModel::computeTidalArguments}).  The BLQ Greenwich phase
lags and the EOT20/FES complex coefficients are both referenced to a
tide argument that carries a \emph{species-dependent Doodson--Warburg
phase bias}.  For the diurnal band the bias is a quarter cycle
($\pm 90^\circ$), because the Cartwright--Tayler--Edden tide-generating
potential refers to a \emph{sine} rather than a cosine when
$(\text{degree}+n_\tau)$ is odd.  The bias per constituent (cycles) is
\begin{equation}
\begin{array}{ccccccccccc}
\text{M2}&\text{S2}&\text{N2}&\text{K2}&\text{K1}&\text{O1}&\text{P1}&\text{Q1}&\text{Mf}&\text{Mm}&\text{Ssa}\\[2pt]
0 & 0 & 0 & 0 & +0.25 & -0.25 & -0.25 & -0.25 & 0 & 0 & 0
\end{array}
\label{eq:sitedisp-dw-bias}
\end{equation}
applied as $\chi_k \mathrel{+}= 2\pi\,(\text{bias}_k)$
(the \texttt{CONSTITUENT\_PHASE\_BIAS\_CYCLES} step of
\texttt{OceanTideModel::computeTidalArguments}; the table itself in
\texttt{OceanTideModel.hpp}).
\textbf{Without this bias every diurnal constituent (K1/O1/P1/Q1) is
wrong in phase by $90^\circ$}, and the whole diurnal OTL contribution is
mismodelled.  The bias must be present, with the signs of
\Cref{eq:sitedisp-dw-bias}.  This settles the question of the diurnal
OTL phase.

\subsubsection{Coefficient sources: SIF, BLQ, and the EOT20 grid}
\label{sec:sitedisp-otl-coeffs}

The per-station amplitudes and phases come from the first available of
four sources (\texttt{Geometry::loadData}, OTL coefficient block): a
Station Information File (SIF), a station BLQ block, the global EOT20
grid, and finally zero.

The SIF gives the best coefficients.  When \texttt{enableStationInfo()}
is set and the SIF has an entry for the RINEX marker, it supplies
per-station three-component FES2014b coefficients.  The SIF stores
Up/East/North amplitude and phase, while the model accumulates in
Up/West/South (\Cref{eq:sitedisp-otl}).  Since
$\text{East}=-\text{West}$ and $\text{South}=-\text{North}$, the
horizontal components keep their amplitude and shift phase by
$+180^\circ$ (\texttt{d.phase\_west = c.east\_phase + 180},
\texttt{d.phase\_south = c.north\_phase + 180}).  Up is unchanged.

Without a SIF entry the coefficients come from a station BLQ block.
Without that, they are synthesised by bilinear interpolation of the
global EOT20 complex load grid, which gives the radial component only.
The grid is sampled at the RINEX header \emph{marker} position
(\texttt{RinexParser::getApproxPositionXyz}), not at the ARP position
state.  The difference is immaterial at this grid resolution, but the
two points are distinct in the code.  With no source the displacement is
zero.

The EOT20 netCDF packs the real
and imaginary load coefficients as \texttt{int16} with a
\texttt{scale\_factor}/\texttt{add\_offset}.  The netCDF C library does
\emph{not} apply these automatically, so the code unpacks them as
$v = v_{\text{int16}}\cdot\text{scale} + \text{offset}$
(\texttt{OceanTideModel::parseEot20Load}: its \texttt{getScaleOffset}
lambda and the manual unpack that follows).  The scale is
$\approx 5.34\times10^{-5}$~m/LSB, and omitting it inflates the
amplitudes by $\sim$20000$\times$.  The complex load
$Z = \mathrm{Re} + i\,\mathrm{Im}$ yields the real signal
$h(t)=\mathrm{Re}\{Z\,e^{i\chi}\} = A\cos(\chi-\phi)$ with
\begin{equation}
A = \sqrt{\mathrm{Re}^2 + \mathrm{Im}^2},\qquad
\phi = \operatorname{atan2}(-\mathrm{Im},\,\mathrm{Re})
\label{eq:sitedisp-otl-blq}
\end{equation}
($\phi$ in BLQ format being the Greenwich phase lag in degrees; same
routine).  Because EOT20 orders its 17
constituents differently from the BLQ convention, the synthesis remaps
the BLQ slot $k$ to the EOT20 index via
$\{4,13,8,3,2,9,10,11,6,7,15\}$
(the \texttt{EOT20\_INDICES\_FOR\_BLQ} table in
\texttt{OceanTideModel::parseEot20Load}).

% ============================================================
\subsection{Pole tide (rotational deformation, IERS 7.1.4)}
\label{sec:sitedisp-pt}

The solid pole tide is the elastic deformation caused by the centrifugal
potential of polar motion.  As the rotation pole moves relative to its
secular mean, the bulge it raises moves with it.  The wobble is the
polar-motion offset from the IERS~2018 secular mean pole.  In
milliarcseconds, with $t_y$ the decimal year,
\begin{equation}
\bar{x}_p = 55.0 + 1.677\,(t_y-2000),\qquad
\bar{y}_p = 320.5 + 3.460\,(t_y-2000)
\label{eq:sitedisp-meanpole}
\end{equation}
(\texttt{GnssCorrections.cpp}, \texttt{computePoleTideCorrection}).  The
post-2018 IGS products used by the engine are referenced to this linear
model.  The wobble components (arcseconds) are
\begin{equation}
\polem{1} = x_p - \bar{x}_p,\qquad
\polem{2} = -\big(y_p - \bar{y}_p\big),
\label{eq:sitedisp-wobble}
\end{equation}
with $x_p,y_p$ the polar-motion values linearly interpolated from the
ERP table at the query epoch
(\texttt{computePoleTideCorrection}; the $\bar{x}_p,\bar{y}_p$
millarcsecond means are scaled to arcseconds in the difference).
\textbf{The wobble must be in arcseconds, not radians.}  The
displacement coefficients below are in metres per arcsecond, so a wobble
in radians scales the whole pole tide by $\sim\!4.85\times10^{-6}$.  The
wobble comment in \texttt{computePoleTideCorrection} records this as a
bug that was once present.

IERS~2010 Eq.~7.24 gives the rotational-deformation displacement in
colatitude $\theta=90^\circ-\geolat$; substituting geodetic latitude
($\sin 2\theta = \sin 2\geolat$, $\cos 2\theta = -\cos 2\geolat$,
$\cos\theta = \sin\geolat$, north $=-S_\theta$) yields the radial,
north, and east components (metres):
\begin{empheq}[box=\fbox]{equation}
\begin{aligned}
\dispPT^{R} &= -0.033\,\sin 2\geolat\,\big(\polem{1}\cos\geolon + \polem{2}\sin\geolon\big),\\
\dispPT^{N} &= -0.009\,\cos 2\geolat\,\big(\polem{1}\cos\geolon + \polem{2}\sin\geolon\big),\\
\dispPT^{E} &= +0.009\,\sin\geolat\,\big(\polem{1}\sin\geolon - \polem{2}\cos\geolon\big),
\end{aligned}
\label{eq:sitedisp-poletide}
\end{empheq}
(\texttt{computePoleTideCorrection}).  \textbf{The east factor is
$\sin\geolat$ ($=\cos\theta$), not $\cos\geolat$.}  A factor of
$\cos\geolat$ is inconsistent with the radial and north terms, which use
geodetic $\geolat$.  It gives a scale error of $\sim\!1.48\times$ at
mid-latitude and the wrong latitude dependence, vanishing at the equator
instead of the pole.  This settles the question of the pole-tide east
term.  The local $(R,N,E)$ vector is rotated into ECEF at the end of
\texttt{computePoleTideCorrection}.

% ============================================================
\subsection{Ocean pole tide (Desai 2002, IERS 7.1.5)}
\label{sec:sitedisp-opt}

Polar motion also redistributes ocean mass.  The load of this ocean pole
tide deforms the crust by a further amount, from below a millimetre to
$\sim$2~mm.  It is added to \Cref{eq:sitedisp-poletide} when the
station's six Desai grid
load coefficients are available
(the \texttt{has\_opt\_coeffs\_} branch of \texttt{computePoleTideCorrection}).  Following Desai~(2002)
\cite{desai2002ocean}, the same wobble $\polem{1},\polem{2}$ drives a
complex pole-tide coefficient
\begin{align}
c^{R} &= \polem{1}\,\gamma_2^{R} + \polem{2}\,\gamma_2^{I},\\
c^{I} &= \polem{2}\,\gamma_2^{R} - \polem{1}\,\gamma_2^{I},
\label{eq:sitedisp-opt-c}
\end{align}
with the IERS/Desai load Love combination
$\gamma_2^{R} = 0.6870,\ \gamma_2^{I} = 0.0036$
(\texttt{kOptGammaReal}/\texttt{kOptGammaImag}).  Each of the six
per-station coefficients $u_c^{R},u_c^{I}$ (north, east, radial real and
imaginary parts), bilinearly interpolated from the Desai grid, gives the
local displacement
\begin{empheq}[box=\fbox]{equation}
\dispOPT^{c} = \kappa\,\big(c^{R}\,u_c^{R} + c^{I}\,u_c^{I}\big),
\qquad c\in\{E,N,U\},
\label{eq:sitedisp-opt}
\end{empheq}
(\texttt{GnssCorrections.cpp}, \texttt{computeOceanPoleTideDisplacement}).
The scale $\kappa$ combines the IERS load-deformation constants with the
arcsecond-to-radian conversion, so the wobble stays in arcseconds as in
\Cref{eq:sitedisp-wobble}:
\begin{equation}
\kappa = \frac{4\pi\,G\,a\,\rho_w\,H_p}{3\,g_E}\,\frac{\pi}{648000},
\qquad
H_p = \sqrt{\tfrac{8\pi}{15}}\,\frac{\Omega^2 a^4}{GM_\oplus},
\label{eq:sitedisp-opt-kappa}
\end{equation}
with $\rho_w = 1025\ \mathrm{kg\,m^{-3}}$,
$a = 6378136.6$~m, $g_E = 9.7803278\ \mathrm{m\,s^{-2}}$,
$\Omega = 7.292115\times10^{-5}\ \mathrm{rad\,s^{-1}}$,
$G = 6.67428\times10^{-11}$, $GM_\oplus = 3.986004418\times10^{14}$
(the \texttt{kOpt*} constants; $\kappa$ assembled in \texttt{oceanPoleTideKappa}).  The whole ocean-pole-tide path
reproduces the NGS M-PAGES implementation and the independent Python
oracle \texttt{scripts/ocean\_pole\_tide\_reference.py} to $\sim\!10^{-12}$
(see the constants' header comment).

% ============================================================
\subsection{Atmospheric pressure loading (S1/S2 thermal tide)}
\label{sec:sitedisp-apl}

Atmospheric pressure loading (APL) is the response of the crust to
changes in surface pressure.  It is $\sim\!2$--$20$~mm and mostly
vertical.  It has two parts.  The $S_1/S_2$ tidal part is
sun-synchronous and thermally driven (diurnal and semidiurnal), and is
predictable from local solar time.  The non-tidal (synoptic) part is
driven by weather systems.  The engine models the \emph{tidal} part
only; the non-tidal part is a gap (see the note below).

The thermal-tide displacement uses per-station diurnal ($S_1$) and
semidiurnal ($S_2$) harmonic coefficients evaluated against the
local-solar-time argument $\Theta_n = n\,(2\pi\,\mathrm{UT}_{\mathrm{frac}}
+ \geolon)$, $n\in\{1,2\}$, with $\mathrm{UT}_{\mathrm{frac}}$ the UT
fraction of day and $\geolon$ the station longitude
(\texttt{GnssCorrections.cpp}, \texttt{computeAtmosphericTideDisplacement}):
\begin{empheq}[box=\fbox]{equation}
\dispAPL^{c} = a^{\cos}_{c,S_1}\cos\Theta_1 + a^{\sin}_{c,S_1}\sin\Theta_1
             + a^{\cos}_{c,S_2}\cos\Theta_2 + a^{\sin}_{c,S_2}\sin\Theta_2,
\qquad c\in\{E,N,U\},
\label{eq:sitedisp-apl}
\end{empheq}
formed in local ENU and rotated to ECEF, with a 1~cm sanity guard.  The
twelve coefficients per station are ordered
$[\,a^{\cos}_E, a^{\sin}_E, a^{\cos}_N, a^{\sin}_N, a^{\cos}_U,
a^{\sin}_U\,]$ for $S_1$ (indices $0$--$5$) then the same six for $S_2$
(indices $6$--$11$).  They are bilinearly interpolated from a
Ray--Ponte-style $S_1/S_2$ load grid
(\texttt{loadAtmosphericTideFromGrid}) at the same RINEX header marker
position as the OTL grid (\Cref{sec:sitedisp-otl-coeffs}).  The
per-epoch displacement is evaluated at the ARP position state.  The term
is gated by \texttt{model\_corrections.enable\_atmospheric\_loading}
(default \texttt{false}).  It reads its grid from
\texttt{files.atm\_tide\_grid\_path}, which defaults to
\texttt{<data>/static/apl/s1\_s2\_load\_cm.dat}.  That filename is
\textsc{provisional} and will be set to the supplied product.  The
coefficients are loaded in \texttt{Geometry::loadData} and projected to
\texttt{apl\_corr} alongside \texttt{otl\_corr} in
\texttt{computeSatelliteGeometry}.  With the flag off, or with no usable
grid, $\dispAPL = \bm{0}$ in \Cref{eq:sitedisp-sum} and runs are
byte-identical.

\paragraph{What is not modelled.}  Only the sun-synchronous $S_1/S_2$
tidal APL is modelled.  The non-tidal (synoptic) pressure loading is
larger and is driven by weather.  Modelling it needs a surface-pressure
field convolved with a Green's function, or a gridded non-tidal loading
product.  It is \emph{not} modelled, and it is the leading candidate for
the next site-displacement term.

% ============================================================
\subsection{Integration into the measurement model}
\label{sec:sitedisp-integration}

Each epoch, the geometry builder evaluates every enabled displacement
and projects it with \Cref{eq:sitedisp-projection}.  The modelled-range
assembly then adds the scalar corrections.  In
\texttt{Geometry::computeSatelliteGeometry} the OTL ECEF displacement is
formed and projected to \texttt{otl\_corr}, the $S_1/S_2$ atmospheric
tide to \texttt{apl\_corr}, and the solid tide is evaluated against the
cached Sun/Moon positions and already returns its LOS-projected scalar.
In \texttt{PppModel} the pole tide (and, when
enabled, the ocean pole tide additively inside it) is projected by
\texttt{computePoleTideProjection} and the base
modelled value (\texttt{computeBaseModeledValue}) sums
\texttt{+ otl\_corr + apl\_corr + solid\_tide}
while the pair processors add
\texttt{+ pole\_tide\_corr}.  The result is the $\geomrange(\rxpos)$ of
\dcref{eq:geom-linear}{the linearised range} with $\rxpos$ replaced by
\Cref{eq:sitedisp-sum}.

The differenced observables of \dcref{sec:bssd}{the BSSD} and
\dcref{sec:dd}{double-difference} formulations contain whatever part of
these corrections survives differencing.  The solid Earth tide and the
pole tide are highly correlated in space and largely cancel over short
baselines.  Ocean loading and atmospheric loading, the latter largely
unmodelled in its non-tidal part, do not cancel.  They are the dominant
differential site-displacement error on long baselines.

\ifSubfilesClassLoaded{\printbibliography}{}
\end{document}
